\section{讨论与结论}
\label{sec:experiments-discussion}

\subsection{结果讨论}

M1 先回答了一个容易忽略的问题：作者代码中的重叠模目标与论文公式中的平方
保真度并不是同一个优化问题。使用五组相同初值时，两种 loss 得到的终态结果也
有明显差别。为了让后续比较采用同一目标，M2--M4 统一使用论文公式对应的
\texttt{paper\_squared}；这一选择并不是因为它在本次实验中的准确率更高。

在这个基础上，M2 分别考察了减少层数、增加量子比特和加入 CZ 三种变化。
1Q-L2、2Q-L2 separable 和 2Q-L2 CZ 都没有在 $\pm0.005$ 范围内保持
1Q-L4 的准确率。CZ 相对 separable 的平均值虽然更高，但五个 seed 中只有两个
出现正差值。M3 将层数、参数量和 level-0 depth 都减少了 $25\%$，却没有通过
准确率判据，优化后的 level-3 depth 也没有下降。这说明在当前设置下，缩小模型
或模板线路并不会自然带来完整流程上的资源收益。

满足实验判据的是 M4。它不改动模型参数，只在推断时调整测量次数；两个模型
都以约 315--330 次平均 shots 接近 fixed 2048 的准确率。这个结果说明有限 shots
模拟中存在节省测量次数的空间，但不是统计等价检验，也不能直接换算为真机运行
时间。由于 M3 没有通过准确率判据，我们也不把两者合并成“剪层与 shots 联合
优化成功”。

\subsection{局限与复现}

本文的正文实验全部在模拟器上完成。训练和确定性评价使用状态向量模拟，M4 从
理想概率中抽取有限 shots，线路也只编译到预先设定的模拟目标。实验没有加入设备
噪声、真机运行或完全相同设置下的经典基线，并且只研究了 circle 数据和五个初始
参数 seed。因此，这些结果不能用来声称硬件加速、量子优势、真机性能或总体统计
显著性。作者设置下的 L8 未收敛结果仍然保留；2Q-L2 CZ 的 level-3 概率误差按
$10^{-10}$ 阈值如实记为未通过，但这不表示结果文件损坏；M3 的失败结果也没有
通过调整准确率阈值来消除。

正文一共包含 45 次优化运行：M1、M2、M3 各 15 次，M4 不重新训练；P0 和 loss
单元测试不计入。汇总检查确认了 45 个互不重复的运行和 246 项实验输入，包括
数据集、索引、源码与测试文件。
每次实验的配置、命令、参数快照、优化状态、\texttt{nfev} 与 SHA-256 均已保存。
最终汇总位于 \path{results/summary/project/20260728T123108.528816Z/}，完整
索引见 \path{PROJECT_EXPERIMENTS.md}。部分早期运行没有同时保存独立的源码
哈希，P0 记录的命令与实际环境中的解释器路径也不完全一致。这些问题不改变已经
保存的数值，但意味着我们不能声称所有运行都能还原到逐字节相同的源码。
Appendix B 是其他组员完成的补充实验，与 M1--M4 分开记录，也不替代正文结果。

\subsection{结论}

本实验表明，比较量子分类器的资源消耗时，损失函数、初始参数、模型参数来源、
编译条件和准确率阈值都需要保持一致。在当前 circle 数据和五个 seed 下，
2Q-L2 的量子比特/层数/CZ 交换和 L4$\to$L3 剪层都没有同时改善准确率与资源；
只有固定模型参数后的 adaptive shots 达到了实验采用的阈值。这里得到的是数值模拟
结果，不能据此证明量子硬件优势。
